\documentclass[10pt,a4paper]{amsart}
\usepackage[margin=19mm]{geometry}
\usepackage{amsmath,amssymb,mathtools}
\usepackage[scr=rsfs,cal=euler]{mathalfa}
\usepackage{xfrac}
\usepackage[unicode,hidelinks,bookmarks=false]{hyperref}
\newcommand{\ie}{i.e.\ }
\newcommand{\cf}{cf.\ }
\newcommand{\resp}{respectively\ }
\newcommand{\Subsection}{\subsection}
\providecommand{\nobreakdash}{\nobreak\hbox{-}}
\setlength{\emergencystretch}{2em}
\multlinegap0pt
\usepackage{mathtools}
\usepackage{xcolor}
\usepackage{amssymb}
\usepackage{enumerate}
\usepackage[normalem]{ulem}
\newtheorem{remark}{Remark}
\newtheorem{lemma}{Lemma}
\newcommand{\K}{\mathbb{K}}
\newcommand{\cck}{C_C (G,\mathbb{K})}
\newcommand{\lck}{L^2 (G,\mathbb{K})}
\newcommand{\n}{\hfill\\}
\title{A Note on the Peter-Weyl Theorem}
\author{Yanga Bavuma, Francesco G. Russo, Elizabeth Stevenson}
\date{Source: arXiv:2603.07105v1 (CC BY 4.0)}
\begin{document}
\maketitle

\noindent\emph{Source and licence.} A Note on the Peter-Weyl Theorem. Version arXiv:2603.07105v1 (CC BY 4.0), distributed by arXiv under the Creative Commons Attribution 4.0 International licence, http://creativecommons.org/licenses/by/4.0/, as recorded in the arXiv licence metadata for this exact version and shown on the abstract page https://arxiv.org/abs/2603.07105v1. Full licence notice: \texttt{licenses/arxiv-2603.07105-CC-BY-4.0.txt}.\par\smallskip

\noindent\textit{Editorial scope.} Counted unit: the article's lemma liftinglemma (source0.tex lines 367--372). The article's complete original proof of it is delivered with it (source0.tex lines 373--406, one \texttt{proof} environment of 34 lines). No mathematics is written, repaired or re-derived: the statement and the proof are the article's own lines, in the article's own notation, and no retained line is substituted or re-flowed.  Retained uncounted as the source-local context the proof needs: lines 344--363.  Nothing was omitted from any retained span, so the package carries no omission record.  No retained line contains a citation, so the wrapper carries no bibliography and no bibliography entry.  No retained line contains a figure environment, an image-inclusion command or drawing code, so no image is delivered and the package declares no asset dependency.  Editorial formatting only, in addition: an AMS \texttt{amsart} preamble that restates the article's own declarations and macros used by the retained lines, namely the environments \texttt{remark}, \texttt{lemma} on separate counters, together with the standard packages those lines need.  The retained environments print in this document as Remark 1, Lemma 1 in order of appearance, whereas the article prints the same items with the numbering of its own full document.  The complete native source of the article as distributed in the arXiv e-print for this exact version is bundled unchanged as \texttt{sources/195919/source0.tex}.
\begin{remark}[Properties of the Lifting Operator]\label{liftingopproperties} Below we will mention a few key properties of the lifting operator that we will need for our deductions later on. Note we show many of these facts only for the $x\in H$ case as the $x\notin H$ case is trivial.
 \begin{enumerate}
 \item\textbf{The Lifting Operator is Linear:}
For all $k_1,k_2\in R(H,\K)$ and $a,b\in\K$ and for all $x\in H$
\[  \ell_G(ak_1+bk_2)(x)=a\ell_G(k_1)(x)+b\ell_G(k_2)(x). \]
 \item\textbf{The Lifting Operator Preserves Positivity:}
 Let $k\in R(H,\K)$ where $k  >0$ for all $x \in H$.  Then
\[ \ell_G(k)(x)=  k(x) \ge 0. \]
 \item\textbf{The Lifting Operator Preserves Shifts:}
 If $k \in R(H,\K)$ and $h\in H$, then for all $x\in H$
\[ \ell_G({}_h k)(x)= \ell_G(k)(xh)=  {}_h \ell_G(k)(x),\]
 therefore the shifting action of $H$ is preserved. %\textcolor{red}{(Questions over this one.)}. 
 \item\textbf{The Lifting Operator Preserves Multiplication:}
 If $k_1,k_2\in R(H,\K)$, then for all $x\in H$ we have
\[ \ell_G(k_1k_2)(x)=\ell_G(k_1)(x)\ell_G(k_2)(x). \] %\textcolor{red}{(Questions over this one.)}
 \item\textbf{The Lifting Operator Preserves the Norm of $\K$:}
 If $k\in R(H,\K)$, then for all $x\in H$
\[ \ell_G(|k|)(x) =|\ell_G(k)(x)|.\]
 \end{enumerate}
\end{remark}

\begin{lemma}[Preservation of Lifting $L^2$-Norms]\label{liftinglemma} Assume that $G$ is a locally compact group with  a nontrivial  compact open  subgroup $H$, and that $\int_G\cdot \;d\lambda$ is a Haar measure on $G$ satisfying (given $\chi_H:G\rightarrow\K$ characteristic function on $H$) $\int_G \chi_H  \; d\lambda=1.$ 
If  $||f||_2$ with $f\in\cck$ is the $L^2$-norm on $\lck$ and  $||k||_{2,H}$ is the $L^2$-norm on $L^2(H,\K)$ with $k\in C(H,\K)$, then  for all $j\in R(H,\K)$ 
 \[
 ||\ell_G(j)||_2=||j||_{2,H}.
 \]
\end{lemma}

\begin{proof}

 First, we claim that the function $\rho:C(H,\K)\rightarrow\K$ given by $\rho (j)\coloneqq \int \ell_G({j})\;d\lambda$ is the normalized Haar measure on $H$. We prove it, checking the definitions and utilizing Remark \ref{liftingopproperties}.  Fix $j,j_1,j_2\in R(H,\K)$ for the below:
 \begin{enumerate}
 \item \textbf{Linearity:} Note for any $a,b\in\K$ we have $\rho(aj_1+bj_2)=\int \ell_G{(aj_1+bj_2)}\;d\lambda=\int a\ell_G{(j_1)}+b\ell_G{(j_2)}\;d\lambda=a\int \ell_G{(j_1)}\;d\lambda+b\int\ell_G{(j_2)}\;d\lambda=a\rho(j_1)+b\rho(j_2)$.\n
 This reasoning uses property (1) of Remark \ref{liftingopproperties}, which follows from the deduction that for all $x\in G$ and $k_1,k_2\in R(H,\K)$ we have 
\[\begin{split}  
 \ell_G(ak_1+bk_2)(x)=&\begin{cases}
                        ak_1(x)+bk_2(x) & \mbox{if} \ x\in H\\
                        0& \mbox{if} \  x\in G\setminus H
                        \end{cases}\\
                        =&a\begin{cases}
                        k_1(x)& \mbox{if} \  x\in H\\
                        0&x\in G\setminus H
                        \end{cases}+b\begin{cases}
                        k_2(x) & \mbox{if} \  x\in H\\
                        0& \mbox{if} \ x\in G\setminus H
                        \end{cases}\\
                    &=a\ell_G(k_1)(x)+b\ell_G(k_2)(x). 
 \end{split}    
\] 
 The properties below of positivity, left invariance, and normalization also rely on the assertions of Remark \ref{liftingopproperties} which may  be  proved in a similar way.
 \item \textbf{Positivity:} Note for $j>0$ we have $\rho(j)=\int \ell_G({j})\;d\lambda$. This is positive as $\ell_G({j})>0$ (Remark \ref{liftingopproperties} property (2)) and Haar integrals are positive.
 
 \item \textbf{Left Invariance:} Note that for $h\in H$ we have $\rho({}_hj)=\int \ell_G{({}_hj)}\;d\lambda=\int {}_h\ell_G({j})\;d\lambda=\int \ell_G({j})\;d\lambda$ (Remark \ref{liftingopproperties} property (3))
 
 \item \textbf{Normalization:} Note that (given \textbf{1}$:H\rightarrow\K$ the constant function on $H$ that takes the value 1 everywhere) $\rho($\textbf{1}$)=\int \ell_G($\textbf{1}$)\;d\lambda=\int \chi_H\;d\lambda=1$
 
 \end{enumerate}
 Next, to show that $||\ell_G({j})||_2=||j||_{2,H}$ we note that (from Remark \ref{liftingopproperties} (4-5))
\[||\ell_G({j})||_2=\sqrt{\int \ell_G({j})\overline{\ell_G({j})}d\lambda}
 =\sqrt{\int |\ell_G({j})|^2d\lambda}
 =\sqrt{\int \ell_G(|j|^2)d\lambda}=\sqrt{\rho(|j|^2)} =\sqrt{\rho(j\;\overline{j})} =||j||_{2,H}.\]
\end{proof}

\end{document}
